\chapter{Results}

% TODO Intro text

% todo significance

\section{Embedding Model Comparison}

\subsection{Embedding Model Selection Criteria}

Embedding models were selected based on a combination of computational feasibility and representational performance. Since embeddings are stored for more than 1.5 million posts, embedding dimensionality directly impacts storage requirements and training efficiency of downstream regression models.

To balance representational capacity with computational constraints, models producing embeddings of approximately 1024 dimensions were preferred. Higher-dimensional embeddings would substantially increase dataset size and model training cost without guaranteed performance gains.

In addition, models with fewer than one billion parameters were prioritized to ensure feasible local execution within available hardware resources. All embedding models were run locally to maintain full control over preprocessing and avoid reliance on external APIs.

Among models meeting these constraints, selection was further guided by reported performance on semantic similarity and classification benchmarks. Since the task involves distinguishing linguistic and stylistic differences between subreddit communities, embedding models with strong semantic discrimination capabilities were considered suitable candidates.

The following section evaluates the selected models empirically within the popularity prediction framework.

\subsection{Evaluation Setup}

To evaluate the impact of embedding model choice on popularity prediction performance, several transformer-based encoder models satisfying the above criteria were compared.

For all experiments in this section, the aggregation strategy was fixed to concatenation of title and selftext embeddings. No dimensionality reduction was applied. All other model components, including the XGBoost regression configuration and training/validation split, were kept constant to ensure a controlled comparison.

Performance metrics are reported as mean values across all subreddits.


\subsection{Results and Interpretation}

\begin{table}[ht]
\centering
\caption{Popularity prediction results for popularity regressor with different encoder models. Results show mean performance across all subreddits.}
\label{tab:embedding_comparison}
\begin{tabular}{lrrrrrr}
\toprule
Features & MAE & RMSE & R² & r & Total \\
\midrule
e5 Large Instruct & 0.1154 & 0.1462 & 0.1680 & 0.4062 & 369,834 \\
Stella 400M & 0.1153 & 0.1459 & 0.1700 & 0.4086 & 369,834 \\
Kalm Mini Instruct & 0.1158 & 0.1465 & 0.1636 & 0.4011 & 369,834 \\
Bilingual & 0.1165 & 0.1473 & 0.1559 & 0.3912 & 369,834 \\
Gwen & 0.1180 & 0.1490 & 0.1361 & 0.3658 & 369,834 \\
\bottomrule
\end{tabular}
\end{table}

Table~\ref{tab:embedding_comparison} reports regression performance across different embedding models. Stella achieves the best overall performance, with the lowest MAE (0.1153), lowest RMSE (0.1459), highest $R^2$ (0.1700), and highest Pearson correlation (0.4086).

The performance differences between Stella and e5 Large Instruct are relatively small, indicating that both models provide strong semantic representations for the popularity prediction task. However, Stella consistently achieves slightly better correlation and variance explained, suggesting improved alignment between predicted and observed upvote ratios.

Kalm Mini Instruct and Bilingual embeddings perform moderately worse, while Gwen exhibits the lowest predictive performance across all reported metrics.

The observed differences across embedding models suggest that higher-quality semantic representations contribute positively to popularity prediction. However, improvements remain moderate, indicating that embeddings alone explain only a limited portion of variance in upvote ratio.

The relatively small gap between the top-performing models also suggests diminishing returns from purely improving semantic encoding without incorporating additional structured features.

Given its consistently superior performance across evaluation metrics, Stella was selected as the baseline embedding model for subsequent experiments. All further aggregation and feature engineering evaluations use Stella embeddings unless otherwise stated.

\subsection{Prompt Conditioning for Stella Embeddings}

To assess whether task-specific prompt conditioning influences embedding quality for popularity prediction, multiple prompt variants were evaluated using the Stella model. All other components of the pipeline were held constant, including aggregation strategy (concatenation), feature configuration (embeddings only), regression model, and training split.

Each prompt variant modifies the instruction context under which embeddings are generated, emphasizing different aspects such as popularity, writing quality, engagement, political leaning, or general content description.

\begin{table}[ht]
\centering
\caption{Popularity prediction results for popularity regressor with different encoder prompts. Results show mean performance across all subreddits.}
\label{tab:prompt_comparison}
\begin{tabular}{lrrrrrr}
\toprule
Features & MAE & RMSE & R² & r & Total \\
\midrule
Stella Popularity & 0.1150 & 0.1456 & 0.1738 & 0.4134 & 369,834 \\
Stella Writing Quality & 0.1148 & 0.1454 & 0.1759 & 0.4158 & 369,834 \\
Stella Engagement & 0.1152 & 0.1459 & 0.1708 & 0.4099 & 369,834 \\
Stella Baseline & 0.1153 & 0.1459 & 0.1700 & 0.4086 & 369,834 \\
Stella Political Leaning & 0.1151 & 0.1456 & 0.1729 & 0.4122 & 369,834 \\
Stella Content & 0.1148 & 0.1454 & 0.1762 & 0.4164 & 369,834 \\
\bottomrule
\end{tabular}
\end{table}

Table~\ref{tab:prompt_comparison} reports the impact of prompt conditioning on regression performance. Differences between prompt variants are relatively small but consistent.

The \textit{Content} prompt achieves the highest performance overall, with the lowest MAE (0.1148), highest $R^2$ (0.1762), and highest Pearson correlation (0.4164). The \textit{Writing Quality} prompt performs nearly identically, while the baseline prompt yields slightly lower performance across all metrics.

Prompts emphasizing popularity or political leaning provide modest improvements over the baseline but do not outperform the more general content-focused variants.

These results suggest that prompt conditioning can influence embedding representations in ways that affect downstream regression performance. However, the magnitude of improvement remains marginal, indicating that embedding quality is relatively robust to prompt phrasing within the tested range.

The superior performance of the content-oriented prompts may indicate that general semantic alignment provides a stronger signal for predicting community approval than narrowly framed task instructions.

Given the negligible performance differences observed across configurations, the Stella model without prompt augmentation continues to serve as the baseline embedding setup for all subsequent aggregation and feature engineering experiments.

\section{Aggregation Strategy Evaluation}

\subsection{Evaluation Setup}

To evaluate the impact of embedding aggregation strategy on popularity prediction performance, multiple methods for combining title and selftext embeddings were compared.

All experiments in this section use the Stella Content embedding configuration. No dimensionality reduction was applied. The XGBoost regression model and train/validation split were kept constant to ensure controlled comparison.

\subsection{Results and Interpretation}

\begin{table}[ht]
\centering
\caption{Popularity prediction results for popularity regressor with different aggregation methods. Results show mean performance across all subreddits.}
\label{tab:aggregation_comparison}
\begin{tabular}{lrrrrrr}
\toprule
Features & MAE & RMSE & R² & r & Total \\
\midrule
Hadamard Product & 0.1211 & 0.1525 & 0.0918 & 0.2843 & 369,834 \\
Mean & 0.1150 & 0.1457 & 0.1731 & 0.4124 & 369,834 \\
Multi-Feature & 0.1153 & 0.1459 & 0.1707 & 0.4094 & 369,834 \\
Conatencation (Baseline) & 0.1153 & 0.1459 & 0.1700 & 0.4086 & 369,834 \\
Min Pooling & 0.1154 & 0.1460 & 0.1691 & 0.4078 & 369,834 \\
Max Pooling & 0.1154 & 0.1461 & 0.1683 & 0.4066 & 369,834 \\
Smart Mean & 0.1152 & 0.1458 & 0.1717 & 0.4107 & 369,834 \\
Title Only & 0.1173 & 0.1483 & 0.1447 & 0.3763 & 369,834 \\
\bottomrule
\end{tabular}
\end{table}

Table~\ref{tab:aggregation_comparison} reports regression performance across different aggregation strategies. Mean pooling achieves the lowest MAE (0.1150) and strong correlation (0.4124), while smart mean performs similarly with slightly improved $R^2$ compared to the baseline concatenation approach.

Concatenation (baseline Stella) achieves competitive performance but does not outperform mean pooling. The multi-feature aggregation, which includes an additional similarity score between title and body embeddings, performs similarly to concatenation but does not provide measurable improvement.

Hadamard product aggregation performs substantially worse than all other methods, with a marked decrease in $R^2$ (0.0918) and correlation (0.2843). This suggests that multiplicative interaction alone is insufficient to preserve predictive information contained in separate title and body embeddings.

Using only the title embedding results in noticeably lower performance compared to combined representations, confirming that body text contributes meaningful signal for popularity prediction.

The performance of smart mean aggregation slightly exceeds that of naive mean pooling and concatenation in terms of explained variance. This supports the hypothesis that conditionally handling body-less posts prevents representational distortion caused by averaging meaningful embeddings with zero vectors.

Although performance differences among pooling-based strategies remain moderate, the results indicate that preserving balanced information from both title and body text is preferable to interaction-only representations.

Given its consistent and stable performance, mean pooling or smart mean aggregation provide practical aggregation strategies with the latter having a minimal dimensionality increase.

\subsection{Effect of Dimensionality Reduction}

To assess whether dimensionality reduction can improve computational efficiency without degrading predictive performance, Principal Component Analysis (PCA) was applied to embedding representations prior to regression.

Both per-subreddit PCA and global PCA (computed across all subreddits) were evaluated. Reduced dimensionalities of 512, 256, and 128 components were tested. All other model components, including aggregation strategy and regression configuration, were kept constant.

\begin{table}[ht]
\centering
\caption{Popularity prediction results for popularity regressor with different dimensionality reductions. Results show mean performance across all subreddits.}
\label{tab:dimensionality_reduction}
\begin{tabular}{lrrrrrr}
\toprule
Features & MAE & RMSE & R² & r & Total \\
\midrule
No Reduction (Baseline) & 0.1153 & 0.1459 & 0.1700 & 0.4086 & 369,834 \\
Per-Subreddit PCA 512 & 0.1153 & 0.1459 & 0.1700 & 0.4086 & 369,834 \\
Global PCA 256 & 0.1156 & 0.1462 & 0.1669 & 0.4049 & 369,834 \\
Global PCA 128 & 0.1154 & 0.1461 & 0.1680 & 0.4064 & 369,834 \\
Global PCA 512 & 0.1207 & 0.1519 & 0.0976 & 0.2630 & 369,834 \\
\bottomrule
\end{tabular}
\end{table}

Table~\ref{tab:dimensionality_reduction} reports regression performance under different dimensionality reduction settings.

Per-subreddit PCA with 512 components yields performance identical to the baseline configuration, indicating that a moderate reduction in dimensionality can preserve predictive signal when applied within individual communities.

Global PCA with 256 and 128 components results in only minor performance degradation, suggesting that substantial dimensionality reduction is possible without dramatically impacting regression quality.

However, global PCA with 512 components exhibits a significant drop in performance, with $R^2$ decreasing to 0.0976 and correlation falling to 0.2630. This suggests that globally computed principal components may fail to preserve community-specific variance patterns critical for predicting popularity.

\section{Feature Contribution Analysis}

\subsection{Temporal Feature Contribution}

To evaluate the contribution of temporal information, time-related features were added individually to the baseline embedding model (Stella embeddings with concatenation). All other model components, including regression configuration and training split, were kept constant.

Each temporal feature group was evaluated independently to isolate its marginal predictive contribution.

\begin{table}[ht]
\centering
\caption{Popularity prediction results for popularity regressor with different author feature combinations. Results show mean performance across all subreddits.}
\label{tab:temporal_features}
\begin{tabular}{lrrrrrr}
\toprule
Features & MAE & RMSE & R² & r & Total \\
\midrule
Full Timestamp Encoding (All Temporal Features) & 0.1151 & 0.1458 & 0.1722 & 0.4115 & 369,834 \\
Hour of Day (Cyclic) & 0.1152 & 0.1459 & 0.1709 & 0.4097 & 369,834 \\
Month of Year (Cyclic) & 0.1152 & 0.1458 & 0.1713 & 0.4104 & 369,834 \\
Time-of-Day Buckets & 0.1152 & 0.1458 & 0.1712 & 0.4099 & 369,834 \\
Day of Week (Cyclic) & 0.1153 & 0.1459 & 0.1701 & 0.4089 & 369,834 \\
Embeddings Only (Baseline) & 0.1153 & 0.1459 & 0.1700 & 0.4086 & 369,834 \\
All Holidays & 0.1154 & 0.1460 & 0.1697 & 0.4081 & 369,834 \\
Weekend Indicator & 0.1153 & 0.1459 & 0.1699 & 0.4085 & 369,834 \\
\bottomrule
\end{tabular}
\end{table}

Table~\ref{tab:temporal_features} reports the effect of individual temporal features on popularity prediction performance.

Adding the complete timestamp encoding yields the strongest improvement among temporal features, increasing $R^2$ from 0.1700 to 0.1722 and Pearson correlation from 0.4086 to 0.4115. Cyclic encodings of hour, month, and time-of-day buckets provide modest but consistent gains over the baseline.

Day-of-week and holiday-related features show minimal impact on predictive performance. In particular, detailed holiday encodings (minor, major, or all holidays) do not improve performance relative to embeddings alone.

Overall, temporal features provide measurable but moderate improvements in predictive performance. This suggests that posting time influences community reception to some extent, but textual features remain the dominant signal for upvote ratio prediction.

% TODO: things like post time might be very relevant for how many likes, views etc post gets, but with our focus on upvote ratio, its doesnt affect how people receive the post itself

The relatively small magnitude of improvement indicates that temporal context acts as a secondary factor rather than a primary driver of community approval.


\subsection{Domain Feature Contribution}

To evaluate the impact of external content source information, domain-based features were added to the baseline embedding model (Stella embeddings with concatenation). 

Three configurations were tested: continuous correlation scores, binary indicators for top and bottom domains, and a combined configuration including all domain-related features. All other model components were kept constant.

\begin{table}[ht]
\centering
\caption{Popularity prediction results for popularity regressor with different author feature combinations. Results show mean performance across all subreddits.}
\label{tab:domain_features}
\begin{tabular}{lrrrrrr}
\toprule
Features & MAE & RMSE & R² & r & Total \\
\midrule
Domain All & 0.1146 & 0.1452 & 0.1792 & 0.4200 & 369,834 \\
Domain Scores & 0.1146 & 0.1452 & 0.1793 & 0.4201 & 369,834 \\
Domain Binary & 0.1153 & 0.1459 & 0.1709 & 0.4097 & 369,834 \\
Embeddings Only (Baseline) & 0.1153 & 0.1459 & 0.1700 & 0.4086 & 369,834 \\
\bottomrule
\end{tabular}
\end{table}

Table~\ref{tab:domain_features} reports the effect of domain-based features on popularity prediction performance.

Adding domain features yields a noticeable improvement over the embedding-only baseline. Both the full domain feature set and the continuous domain score configuration increase $R^2$ from 0.1700 to approximately 0.1793 and improve Pearson correlation from 0.4086 to 0.4201.

In contrast, the binary domain indicators alone do not provide meaningful improvement over the baseline.

The strong performance of continuous domain score features suggests that subreddit communities exhibit systematic differences in how linked external content is received. However, domain features should not be interpreted purely as reputation signals. 

For link posts, only the title text is embedded, while the full article content remains unobserved by the model. Domain information may therefore act as a proxy for latent characteristics of the linked articles, such as writing style, topical framing, or ideological alignment, which are not captured by title embeddings alone.


\subsection{Sentiment Feature Contribution}

To evaluate whether affective signals contribute to popularity prediction, sentiment features derived from VADER were added to the baseline embedding model (Stella embeddings with concatenation). 

Multiple configurations were tested, varying both the text source (title, body, combined, or max-pooled) and the type of sentiment representation (raw scores, binary indicator, categorical encoding, or combined features).

\begin{table}[ht]
\centering
\caption{Popularity prediction performance with different sentiment feature configurations. Baseline corresponds to Stella embeddings with concatenation and no additional structured features. Results show mean performance across all subreddits.}
\label{tab:sentiment_features}
\begin{tabular}{lrrrrrr}
\toprule
Features & MAE & RMSE & R² & r & Total \\
\midrule
Title & Body Separate (Vader Scores) & 0.1152 & 0.1459 & 0.1709 & 0.4101 & 369,834 \\
Title & Body Separate (All Sentiment Features) & 0.1152 & 0.1459 & 0.1710 & 0.4101 & 369,834 \\
Title & Body Separate (Binary Sentiment) & 0.1153 & 0.1459 & 0.1700 & 0.4086 & 369,834 \\
Embeddings Only (Baseline) & 0.1153 & 0.1459 & 0.1700 & 0.4086 & 369,834 \\
Max-Pooled (All Sentiment Features) & 0.1153 & 0.1459 & 0.1708 & 0.4097 & 369,834 \\
Combined Text (All Sentiment Features) & 0.1153 & 0.1459 & 0.1700 & 0.4088 & 369,834 \\
\bottomrule
\end{tabular}
\end{table}

Table~\ref{tab:sentiment_features} reports the effect of sentiment features on prediction performance.

Across configurations, sentiment features provide only marginal improvements over the embedding-only baseline. The strongest configuration (separate sentiment scores for title and body) increases $R^2$ from 0.1700 to approximately 0.1710 and Pearson correlation from 0.4086 to roughly 0.4101.

Other configurations yield nearly identical performance to the baseline.

These results suggest that coarse-grained sentiment signals contribute limited additional predictive value beyond transformer-based embeddings. Since modern embeddings already encode contextual and affective information implicitly, explicit sentiment scores appear largely redundant in this setting.

The relatively small magnitude of improvement indicates that community approval is not strongly determined by simple positive or negative sentiment alone, but rather by more complex linguistic, contextual, and structural factors.

\subsection{TF-IDF Feature Contribution}

To assess whether lexical frequency-based signals provide additional predictive value beyond transformer embeddings, TF-IDF-derived features were added to the baseline embedding model (Stella embeddings with concatenation).

Three categories were evaluated: coverage features, correlation-based TF-IDF scores, and binary indicators for high-impact terms. A combined configuration including all TF-IDF features was also tested. All other model components were kept constant.

\begin{table}[ht]
\centering
\caption{Popularity prediction performance with TF-IDF-derived feature additions. Baseline corresponds to Stella embeddings with concatenation and no additional structured features. Results show mean performance across all subreddits.}
\label{tab:tfidf_features}
\begin{tabular}{lrrrrrr}
\toprule
Features & MAE & RMSE & R² & r & Total \\
\midrule
TF-IDF All Features & 0.1153 & 0.1460 & 0.1695 & 0.4081 & 369,834 \\
TF-IDF Binary & 0.1153 & 0.1460 & 0.1697 & 0.4081 & 369,834 \\
TF-IDF Scores & 0.1153 & 0.1459 & 0.1700 & 0.4086 & 369,834 \\
Embeddings Only & 0.1153 & 0.1459 & 0.1700 & 0.4086 & 369,834 \\
TF-IDF Coverage & 0.1153 & 0.1459 & 0.1700 & 0.4086 & 369,834 \\
\bottomrule
\end{tabular}
\end{table}

Table~\ref{tab:tfidf_features} reports the effect of TF-IDF-derived features on regression performance.

Across all configurations, TF-IDF features do not yield measurable improvement over the embedding-only baseline. Performance metrics remain effectively unchanged, with $R^2$ and Pearson correlation identical to or slightly below baseline values.

Even the correlation-based TF-IDF score features, which encode term-level relationships with popularity, fail to provide additional predictive signal beyond what is already captured by transformer embeddings.

These findings suggest that dense transformer embeddings already encode the majority of useful lexical information relevant for predicting upvote ratio. The explicit inclusion of TF-IDF-based lexical statistics appears largely redundant in this setting.

This result highlights the representational strength of contextual embeddings compared to traditional bag-of-words approaches, particularly when large-scale pretrained models are used.


\subsection{Author Feature Contribution}

To evaluate the impact of author-level information, author metadata features were added to the baseline embedding model. Due to missing or inaccessible author accounts (e.g., deleted or unavailable users), complete author metadata was not available for all posts.

As a result, experiments involving author features were conducted on a reduced subset of 220,855 posts. Approximately 40\% of posts were excluded because required author metadata could not be retrieved.

To ensure fair comparison, the baseline model reported in this section (``Embeddings Only – Author Filtered'') was evaluated on the same reduced subset without incorporating author features.

\begin{table}[ht]
\centering
\caption{Popularity prediction performance with author-related feature additions. Results are reported on the subset of posts for which complete author metadata was available. Baseline corresponds to embeddings only on the same filtered subset.}
\label{tab:author_features}
\begin{tabular}{lrrrrrr}
\toprule
Features & MAE & RMSE & R² & r & Total \\
\midrule
All Author Features & 0.1090 & 0.1390 & 0.2016 & 0.4444 & 220,855 \\
Account Age & 0.1108 & 0.1412 & 0.1782 & 0.4180 & 220,855 \\
Link Karma & 0.1113 & 0.1415 & 0.1737 & 0.4128 & 220,855 \\
Comment Karma & 0.1113 & 0.1414 & 0.1748 & 0.4141 & 220,855 \\
Total Karma & 0.1112 & 0.1413 & 0.1758 & 0.4152 & 220,855 \\
Moderator Status & 0.1126 & 0.1429 & 0.1576 & 0.3931 & 220,855 \\
Embeddings Only (Author-Available Subset) & 0.1126 & 0.1429 & 0.1571 & 0.3923 & 220,855 \\
Email Verified & 0.1125 & 0.1428 & 0.1594 & 0.3955 & 220,855 \\
Premium (Gold) Status & 0.1126 & 0.1429 & 0.1573 & 0.3928 & 220,855 \\
\bottomrule
\end{tabular}
\end{table}

Table~\ref{tab:author_features} reports regression performance with individual and combined author features.

Incorporating all author-related features yields a substantial improvement in predictive performance, increasing $R^2$ from 0.1571 (baseline on filtered subset) to 0.2016 and Pearson correlation from 0.3923 to 0.4444.

Among individual features, account age and total karma provide the strongest gains, while moderator status and premium (Gold) status contribute comparatively little predictive value.

The magnitude of improvement indicates that author-related variables are strongly associated with post reception. However, this association should not be interpreted as evidence that users directly reward reputation or account status.

At least two plausible explanations exist. First, more experienced users—reflected in higher account age or accumulated karma—may better understand subreddit norms and audience expectations, resulting in higher-quality or better-aligned submissions. Second, reputation-related signals may influence how posts are perceived within the community, independent of intrinsic content quality.

The present analysis cannot distinguish between these explanations. It demonstrates only that author-level characteristics are statistically associated with upvote ratio and improve predictive performance within the regression framework relative to embeddings alone on the same data subset.

Because approximately 40\% of posts lacked retrievable author metadata due to account deletion or archival limitations, the author-based experiments were conducted on a reduced sample. Direct comparison to full-dataset experiments is therefore not appropriate. However, the embedding-only baseline trained on the identical filtered subset provides an internally consistent reference, ensuring that the observed performance differences are attributable to the inclusion of author features rather than changes in sample composition.

\paragraph{Feature Contribution Summary.}

Taken together, the preceding feature analyses reveal distinct differences in predictive contribution across feature groups.

Overall, author-level and domain-based features provide the largest predictive gains, while temporal effects are moderate and sentiment and TF-IDF features contribute minimally beyond embeddings.

\section{Combined Feature Configurations}

This section evaluates regression performance when multiple feature groups are combined. The goal is to determine whether predictive gains observed for individual feature groups accumulate when integrated into a single model.

Results will be reported separately for:
(i) the full dataset without author features, and 
(ii) the author-available subset including user-level variables.


\section{Cross-Subreddit Generalization}

To evaluate the extent to which popularity dynamics transfer across communities, three alternative training strategies were compared:

(1) \textit{Subreddit-specific models}, where a separate regression model is trained for each subreddit.

(2) \textit{Leaning-group models}, where one model is trained per political leaning category.

(3) \textit{Combined model}, where a single model is trained on all subreddits jointly.

In addition, a variant of the combined model was evaluated that includes an explicit subreddit indicator feature.

All configurations use identical embedding and regression settings to ensure comparability.


\begin{table}[ht]
\centering
\caption{Regression performance by model grouping strategy. MAE and $r$ are averaged across subreddits. For the leaning-group models, only the subreddits matching each leaning category are included before averaging.}
\label{tab:model_grouping}
\begin{tabular}{lrrr}
\toprule
Model Type & MAE & $R^2$ & $r$ \\
\midrule
Subreddit-specific  & 0.1153 & 0.1700 & 0.4086 \\
Leaning-group       & 0.1217 & 0.1087 & 0.3540 \\
Combined + subreddit feature & 0.1219 & 0.0984 & 0.3270 \\
Combined            & 0.1298 & -0.0027 & 0.2887 \\
\bottomrule
\end{tabular}
\end{table}

Table~\ref{tab:model_grouping} reports average performance across subreddits for each training strategy.

Subreddit-specific models achieve the strongest performance, with mean Pearson correlation of 0.4086 and $R^2$ of 0.1700. Leaning-group models exhibit a noticeable performance drop, while the fully combined model performs substantially worse, with mean $R^2$ approaching zero.

Introducing an explicit subreddit feature into the combined model partially recovers performance but remains clearly below subreddit-specific training.

These results indicate that popularity dynamics are strongly community-dependent. Training a global model across heterogeneous subreddits introduces substantial noise that reduces predictive accuracy. Even grouping subreddits by political leaning is insufficient to fully capture community-specific variation.

The partial recovery achieved by adding a subreddit indicator feature suggests that shared structural patterns exist across communities, but these patterns require explicit conditioning on community identity to be effectively leveraged.

\clearpage

\begin{longtable}{lrrrr}
\caption{Per-subreddit Pearson $r$ by training strategy. \textit{Specific} trains a separate model per subreddit; \textit{Leaning} trains one model per political leaning group; \textit{Combined} trains a single model on all subreddits; \textit{Combined + Sub.} adds a subreddit indicator feature to the combined model.}
\label{tab:per_subreddit_pearson} \\
\toprule
Subreddit & Specific & Leaning & Combined & Combined + Sub. \\
\midrule
\endfirsthead
\multicolumn{5}{l}{\small\textit{(continued from previous page)}} \\
\toprule
Subreddit & Specific & Leaning & Combined & Combined + Sub. \\
\midrule
\endhead
\midrule
\multicolumn{5}{r}{\small\textit{(continued on next page)}} \\
\endfoot
\bottomrule
\endlastfoot
\multicolumn{5}{l}{\textit{Right-leaning}} \\
\quad Conservative       & 0.3760 & 0.2730 & 0.2388 & 0.3045 \\
\quad TrueChristian      & 0.4288 & 0.3687 & 0.3685 & 0.3722 \\
\quad Libertarian        & 0.5032 & 0.3941 & 0.3587 & 0.4173 \\
\quad mensrights         & 0.4189 & 0.3527 & 0.3378 & 0.3435 \\
\quad progun             & 0.4401 & 0.3358 & 0.3357 & 0.3605 \\
\midrule
\multicolumn{5}{l}{\textit{Left-leaning}} \\
\quad Liberal            & 0.4129 & 0.3665 & 0.2648 & 0.3460 \\
\quad democrats          & 0.4298 & 0.4033 & 0.2372 & 0.3231 \\
\quad racism             & 0.3441 & 0.3060 & 0.2708 & 0.2668 \\
\quad anarchism          & 0.4082 & 0.3919 & 0.2981 & 0.3450 \\
\quad TrueAtheism        & 0.3928 & 0.3494 & 0.2768 & 0.3345 \\
\midrule
\multicolumn{5}{l}{\textit{Neutral}} \\
\quad PoliticalDiscussion & 0.5316 & 0.5049 & 0.3837 & 0.4679 \\
\quad changemyview        & 0.3604 & 0.3091 & 0.2001 & 0.2285 \\
\quad geopolitics         & 0.3417 & 0.2762 & 0.2401 & 0.2574 \\
\quad DebateReligion      & 0.4568 & 0.4125 & 0.2766 & 0.2812 \\
\midrule
\multicolumn{5}{l}{\textit{Apolitical}} \\
\quad DecidingToBeBetter  & 0.3445 & 0.2975 & 0.2886 & 0.2979 \\
\quad truegaming          & 0.4699 & 0.3699 & 0.2813 & 0.3629 \\
\quad Fantasy             & 0.4756 & 0.4411 & 0.3806 & 0.4011 \\
\quad Frugal              & 0.3269 & 0.2710 & 0.2149 & 0.2498 \\
\quad Coffee              & 0.3018 & 0.2817 & 0.2329 & 0.2522 \\
% already included in table before this one
% \midrule
% \textbf{Mean}             & \textbf{0.4086} & \textbf{0.3540} & \textbf{0.2887} & \textbf{0.3270} \\
\end{longtable}

Table~\ref{tab:per_subreddit_pearson} provides a per-subreddit comparison of Pearson correlation across training strategies.

Across nearly all communities, subreddit-specific models achieve the strongest performance. The performance gap is particularly pronounced for politically active subreddits such as PoliticalDiscussion, Libertarian, and democrats.

Leaning-group models typically outperform the fully combined model, indicating that grouping by ideological orientation preserves some transferable structure. However, leaning-level training remains consistently below subreddit-specific performance. Adding a subreddit indicator feature to the combined model improves results across most communities but does not close the gap.

These findings suggest that while some shared engagement patterns exist across ideologically similar communities, substantial variation persists at the individual subreddit level. Popularity dynamics therefore appear to be shaped by both broader ideological alignment and localized community norms.

\subsection{Leaning-Level Transferability}

\begin{table}[ht]
\centering
\caption{Cross-leaning generalisation: mean Pearson $r$ when a leaning-group model is evaluated on subreddits from each leaning category. Diagonal (bold) shows within-group performance; off-diagonal shows cross-group transfer.}
\label{tab:cross_leaning_pearson}
\begin{tabular}{lrrrr}
\toprule
\multirow{2}{*}{Train Model} & \multicolumn{4}{c}{Evaluated on} \\
\cmidrule(lr){2-5}
 & Right & Left & Neutral & Apolitical \\
\midrule
Right model      & \textbf{0.3447} & 0.2180 & 0.1699 & 0.1571 \\
Left model       & 0.2479 & \textbf{0.3604} & 0.1335 & 0.1478 \\
Neutral model    & 0.1425 & 0.1713 & \textbf{0.3750} & 0.1211 \\
Apolitical model & 0.1862 & 0.1625 & 0.1449 & \textbf{0.3308} \\
\bottomrule
\end{tabular}
\end{table}

Table~\ref{tab:cross_leaning_pearson} reports cross-leaning generalisation performance. Each cell shows the mean Pearson correlation when a leaning-group model is evaluated on subreddits belonging to a particular ideological category. Diagonal entries represent within-group performance, while off-diagonal entries capture cross-group transfer.

A clear diagonal structure is visible. Models trained within a given ideological category perform substantially better on subreddits of the same category than on those of different leanings. For example, the right-leaning model achieves a mean $r$ of 0.3447 on right-leaning subreddits but drops markedly when evaluated on left-leaning (0.2180), neutral (0.1699), or apolitical (0.1571) communities.

This pattern holds consistently across all leaning categories, indicating that engagement dynamics exhibit partial ideological clustering. However, even within-group performance remains below subreddit-specific training, demonstrating that ideological similarity alone does not fully account for community-specific popularity patterns.

Taken together, these results suggest a hierarchical structure of generalizability: subreddit-specific models perform best, leaning-group models capture intermediate transferable structure, and fully combined models without explicit community conditioning perform worst.


\section{Regression vs Classification}

To evaluate whether directly modeling popularity as a binary classification task provides advantages over regression, a classification model was trained to predict whether a post belongs to the “popular” class. The popularity threshold was defined using the median upvote ratio of the training split.

For comparability, regression predictions were also converted into binary labels using the same training median threshold. Performance was evaluated on the validation split using standard classification metrics.

Because the median threshold was computed on the training split only, the validation split does not contain exactly 50\% popular posts. As a result, the positive class proportion differs slightly from 50\% in evaluation.

\begin{table}[ht] 
    \centering 
    \caption{Popularity prediction performance using classification. Baseline corresponds to embeddings only on the same filtered subset.} 
    \label{tab:classification} 
    \begin{tabular}{lrrrrrr} 
        \toprule Features & Precision & Recall & F1-Macro & Accuracy & Total & \% Popular \\ \midrule Baseline Classification & 0.6841 & 0.6223 & 0.6513 & 0.6341 & 369,834 & 50.00\% \\ Baseline Regression as Classification & 0.6369 & 0.6718 & 0.6341 & 0.6356 & 369,834 & 54.05\% \\ 
        \bottomrule 
    \end{tabular} 
\end{table}

Table~\ref{tab:classification} compares direct classification with regression-based thresholding. Precision and recall values refer to the positive (“popular”) class unless otherwise specified. Macro F1-score is reported to account for both classes symmetrically.

Direct classification achieves a higher macro F1-score (0.6513) compared to regression-derived classification (0.6341). However, regression-based classification yields higher recall and comparable overall accuracy, while predicting a larger share of posts as popular (54.05\%), reflecting the transfer of the training median threshold to a slightly shifted validation distribution.

Importantly, the deployment objective of this thesis is asymmetric. In the live experiment, posts predicted as popular are actively published, whereas posts predicted as unpopular are discarded and can be replaced with newly generated candidates. Consequently, false positives (predicting a post as popular when it is not) are more costly than false negatives.

From this perspective, precision on the positive class is more critical than recall. Although direct classification achieves marginally better F1-score, regression preserves a continuous prediction scale that enables flexible threshold adjustment to prioritize precision when required.

Because regression supports ranking and precision-oriented filtering without retraining, it provides a more suitable framework for the generation-and-selection pipeline used in the live experiment.


\section{Hyperparameter Optimization}

Initial experiments were conducted using a fixed XGBoost configuration to enable controlled comparison of embedding and feature effects. After identifying promising model configurations, hyperparameter optimization was performed to assess potential additional performance gains.


\section{Final Model Configuration}

\section{Live Experiment Evaluation}